\documentclass{article}
\usepackage{amsmath}
\usepackage{mathtools}
\usepackage{amsfonts}
\usepackage{parskip}
\setlength\parindent{0pt}

\title{Career Development Plan: Layer 1 - Strategic Foundation - DRAFT \\ [1em]
	\large Jonathan Erdmann \\ [1em]
}

\begin{document}

\maketitle

This document is the strategic foundation of my Career Development Plan and serves as the anchor for all planning activity that 
follows. Content in this layer operates at the level of enduring truth.
Statements here should remain stable across the full planning horizon and require deliberate annual review before they are changed.
Every element in Layers 2 and 3 must be traceable to something written here.

\section{Personal Mission}

To apply deep domain expertise and rigorous analytical thinking to complex, high-stakes problems in aerospace and defense, 
owning the full arc from problem definition to decision, such that consequential decisions can be directly traced to the 
quality of my work.

\section{Positioning Thesis}

Two decades in aerospace and defense across structural integrity programs, mission-critical systems, engineering organizations,
and workforce management at scale has produced deep domain fluency, technical judgment, and a firsthand understanding of how
programs get executed and how organizations perform. An Executive MBA in finance and strategy adds to the quantitative and analytical
toolkit to operate in corporate strategy, corporate development, and investment functions. Together these form a professional profile
that can move fluidly between technical substance, organizational dynamics, and strategic implications. These capabilities 
are a unique and valuable combination and are consistently relevant in aerospace and defense.

\section{Target Functions}

%Target functions are sequenced deliberately. Near-term targets inform the operating model immediately. Medium and long-term targets
%inform capability development decisions and are actively revisited at each annual review.

\subsection{Near-Term Primary Target: Corporate Strategy}

The internal advisory function supporting C-suite decision-making at large firms. Owns strategic planning cycles, competitive
intelligence, market sizing, portfolio tradeoffs, and build/buy/partner analysis. Output is generally recommendations and 
frameworks, not transactions or operational decisions. The work is analytical and influence-driven, shaping decisions without 
always owning their execution. At a large prime this function sits close to the CEO and business unit presidents, making 
it one of the highest-visibility staff roles available to someone without a P\&L.

\textbf{Entry Angle:} The positioning thesis directly supports this function --- domain fluency, technical judgment, and quantitative
finance training are precisely differentiators of a strong corporate strategy candidate at a defense prime from a 
generalist consultant. The gap to close is the strategy title itself and demonstrated exposure to strategic planning processes, 
which the operating model must address explicitly.

\subsection{Medium-Term Target: Corporate Development}

This is the transaction-execution function: Owns M\&A sourcing, financial modeling, due diligence, deal structuring, 
and integration planning. Success is measured in closed transactions and not recommendations to decision makers.
At a large prime this function operates with significant financial authority and works directly with executive 
leadership and external advisors on inorganic growth decisions.

\textbf{Entry Angle:} Direct entry without prior corporate development or investment banking experience is uncommon. 
The credible path is through corporate strategy first, where exposure to build/buy/partner analysis, 
portfolio reviews, and strategic transactions creates a natural bridge. The operating model should include 
deliberate M\&A adjacent activities --- whether internal transaction support, partnership structuring, 
or strategic initiative ownership --- that begin building this credential during Phase 2.

\newpage
\subsection{Long-Term Option: Corporate Venture Capital} 

The investment function focused on minority stakes in early-to-growth stage companies, typically housed within 
a prime's dedicated investment arm. Balances strategic fit activities like technology scouting, capability 
adjacencies, ecosystem development with financial return. In aerospace and defense, CVC often functions 
as a forward-deployed R\&D pipeline, identifying technologies before they are mature enough for acquisition. 
Investment theses are evaluated on both strategic relevance and commercial viability.

\textbf{Entry Angle:} Technical domain depth in aerospace and defense systems is a genuine differentiator in this function --- the
ability to evaluate whether a technology is real, a team is credible, and a capability is strategically relevant may be
rare among people with the business credentials to operate in an investment role. The credible path runs through 
corporate development, where deal evaluation and financial modeling skills are formalized, or through sufficient 
ecosystem visibility that a CVC role becomes a direct lateral. 
Preserved as optionality and not actively targeted until Phase 2 is complete.

\newpage

\section{Known Gaps}

Gaps are assessed across two domains --- competency and positioning --- and three time horizons corresponding 
to the target function sequence. Near-term gaps are assessed in detail. Medium and long-term gaps are 
assessed directionally and revisited at each annual review.

Competency gaps address the ability to perform and thrive once in the role. Positioning gaps address the ability to be considered for
and selected into the role. Both must be addressed. Closing one at the expense of the other produces either an unprepared candidate
or an invisible one.

Competency is assessed first because execution of the mission depends on the quality of the work, not the quality of the positioning.

\subsection{Near Term: Corporate Strategy}

\subsubsection{Competency Gaps}

\textbf{Strategic Analysis \& Frameworks}
Some exposure through coursework and work experience but not yet practiced independently at the level required to structure and solve
ambiguous, high-stakes problems without guidance. Formal frameworks have been introduced through the EMBA but applied experience in a
corporate strategy context is non-existent. This is the most critical competency gap to close given the centrality of 
structured analytical thinking to the target function.

\textbf{Financial Acumen}
Foundational understanding established through EMBA coursework but limited applied experience building, stress-testing, or presenting
financial models in a business context. Fluency with financial statements and capital allocation logic exists at a conceptual level but
has not yet been exercised in a corporate strategy or corporate development setting.

\textbf{Industry \& Domain Knowledge}
Deep operational knowledge from within the industry but limited breadth across the competitive landscape, program economics at the
enterprise level, and technology trends outside areas of direct experience. The domain depth is a genuine asset; the gap is in breadth
and in the ability to synthesize across the landscape rather than within a specific program or system.

\textbf{Executive Communication}
Capable of structuring and delivering findings to senior leadership with significant room to grow at more senior levels. 
The gap is in consistency, precision under scrutiny, and the ability to compress complex analysis into the format and 
cadence that corporate strategy demands. Significant gap when communicating in settings where dominant voices frequently
insert themselves.

\textbf{Stakeholder Navigation}
Capable of navigating organizational complexity and building credibility across functions without direct authority. The gap is in
operating at the corporate level where political complexity, ambiguity of authority, and the stakes of decisions are 
materially higher.

\subsubsection{Positioning Gaps}

\textbf{Professional Narrative}
A narrative exists grounded in two decades of A\&D experience but is not yet coherent or tailored to corporate strategy. Current
profile and resume position as an engineering manager, not a strategist. The work required is reframing existing experience ---
which includes genuine strategic work products --- into a story that is legible and compelling to a decision-maker in the target function.

\textbf{External Visibility}
A LinkedIn profile exists but is not yet visible to the right people in the target function. No presence in corporate strategy or
corporate development circles at large primes. The gap is both reach and relevance.

\textbf{Network}
Broad professional network but not targeted to corporate strategy or corporate development functions at large primes. Limited
relationships with people who could surface my name when an opportunity arises or provide insider context on how these functions operate.

\textbf{Internal Positioning}
Some visibility to senior leadership at Collins through current role performance but not yet associated with strategic capability or
ambition beyond engineering management. No executive sponsor actively advocating for development or advancement into a strategy-adjacent role.

\textbf{Credentials}
EMBA in progress; completion expected May 2027. No strategy-adjacent title or role held to date. Like professional 
narrative, part of this gap is due to an engineering-first background.

\subsection{Medium-Term: Corporate Development}

To be assessed directionally following completion of near-term gap closure. Key anticipated gaps include transaction experience,
financial modeling depth, and M\&A process knowledge. Full assessment to be conducted at the Phase 2 annual review.

\subsection{Long-Term: Corporate Venture Capital}

To be assessed at a high level following completion of medium-term gap closure. Key anticipated gaps include investment thesis
development, startup evaluation, and portfolio management experience. Full assessment to be conducted at the Phase 3 annual review.

\newpage

\section{Breakthrough Goals}

\textbf{Goal 1: Complete the SMU Executive MBA}

This is the most time-bound and definitive goal in the plan. Completion closes the credential gap and activates the
quantitative finance and strategic analysis toolkit required to perform in the target function. Graduating with distinction 
signals the seriousness of the credential to a hiring decision-maker although is not necessary and is harder given difficulties
in Term 3 due to Achilles tendon rupture.

\textbf{Goal 2: Establish a coherent, credible professional narrative tailored to corporate strategy and a targeted network
of relevant relationships} 

Success condition: the right people inside Collins and within the broader A\&D corporate strategy community know who I am, understand
what I am building toward, and would surface my name when a relevant opportunity arises. Internal positioning takes priority in the
near term given the long-term Collins-first strategy.

\textbf{Goal 3: Transition into a strategy-adjacent role or formal corporate strategy function within Collins Aerospace
at the Senior Manager level or above}

Collins is the first-move target because existing relationships, organizational knowledge, and demonstrated performance reduce the
credibility gap that an external employer would require evidence to close. An external move to a corporate strategy function at
another company within RTX is the contingency if internal advancement stalls, assessed at the annual review.

\textbf{Goal 4: Own at least one consequential strategic work product within Collins Aerospace where the output directly
informed a senior executive decision.}

Success condition: a market assessment, portfolio recommendation, competitive analysis, or strategic initiative with a traceable
outcome where the link between the work and the decision is direct and attributable.

\textbf{Goal 5: Be positioned credibly for a corporate development role at a large A\&D prime, with M\&A adjacent exposure
and financial modeling depth sufficient to compete as a serious candidate.}

This goal preserves the medium-term target without conflating it with the near-term transition. 
The goal is achieved through the combination of corporate strategy credentialing, deliberate pursuit of M\&A adjacent 
activities, and deepening financial acumen during Phases 2 and 3.

\end{document}
